% ===============================================================
% 总结与展望
% ===============================================================
\section{总结与展望}

% ---------------------------------------------------------------
\begin{frame}{总结与展望}
\begin{columns}[T,onlytextwidth]
\begin{column}{0.57\textwidth}
\textbf{工作总结}

\vspace{1mm}
\begin{tcolorbox}[rustcard,title={面向框内核的覆盖制导灰盒模糊测试}]
  插桩边界与 TCB 边界重合；用 Rust 在 OSTD 中实现无递归、无死锁、
  抵御敌意共享内存的覆盖收集；提供 Linux 兼容的 \code{/dev/kcov}，
  并让 \syz{} 识别 Rust panic
\end{tcolorbox}
\vspace{1mm}
\begin{tcolorbox}[syzcard,title={自适应内核生命周期管理}]
  带折扣的 LinUCB 在继续、检查点、软重启、硬重启之间决策，降低整机重启开销
\end{tcolorbox}
\vspace{1mm}
\begin{tcolorbox}[tcbcard,title={完整的工程闭环}]
  从打补丁到统计评估的一站式流程，支持 \aster{} v0.18.1 与 v0.19.0
\end{tcolorbox}
\end{column}
\begin{column}{0.39\textwidth}
\textbf{未来工作}

\vspace{1mm}
\begin{itemize}\small
  \item \textbf{多架构}：从 x86\_64 扩展到 ARM64、RISC-V
  \item \textbf{异步覆盖}：面向 \code{async}/\code{await} 状态机的细粒度反馈
  \item \textbf{规约生成}：从 AST 与类型系统中提取 syzlang 规约
\end{itemize}

\vspace{2mm}
\centering
\qrcode[height=1.9cm]{\repourl}

\vspace{1mm}
{\scriptsize\faGithub\ \url{\repourl}}
\end{column}
\end{columns}
\end{frame}

% ---------------------------------------------------------------
% 致谢页使用白色背景（metropolis 的 standout 取 palette primary 配色）
{\setbeamercolor{palette primary}{fg=SAText,bg=white}
\begin{frame}[standout]
  \vspace*{0pt}
  {\LARGE 谢谢！}

  \vspace{3mm}
  {\large 敬请各位评委批评指正}

  \vspace{5mm}
  \begin{tcolorbox}[enhanced,colback=white,colframe=SATcbFill,arc=3pt,center,
      width=2.9cm,halign=flush center,left=2pt,right=2pt,top=2pt,bottom=2pt]
    \qrcode[height=2.4cm]{\repourl}
  \end{tcolorbox}

  {\scriptsize\url{\repourl}}
\end{frame}
}
